% !TeX program = xelatex
% 从 docs/latex/ai-authoring 运行 xelatex design.tex；成品位于 output/pdf/Redflag-AI-Native-Design.pdf。
\documentclass[UTF8,11pt,fontset=fandol,a4paper]{ctexart}
\usepackage[margin=23mm,top=24mm,bottom=21mm]{geometry}
\usepackage{tikz,tabularx,array,booktabs,fancyhdr,hyperref}
\usetikzlibrary{arrows.meta,positioning,fit,backgrounds}
\definecolor{ink}{HTML}{19363D}\definecolor{teal}{HTML}{176C65}\definecolor{muted}{HTML}{526B70}
\hypersetup{colorlinks=true,urlcolor=teal,linkcolor=teal,pdftitle={面向人的 Redflag AI 原生编辑架构}}
\ctexset{section={format=\Large\bfseries\color{ink},afterindent=true},subsection={format=\large\bfseries\color{teal},afterindent=true}}
\setlength{\parindent}{2\ccwd}\setlength{\parskip}{.55\baselineskip}
\pagestyle{fancy}\fancyhf{}\lhead{\small Redflag\enspace / \enspace 引擎设计}\rhead{\small 2026.11.22}\rfoot{\small \thepage}
\newcommand{\src}[2]{\href{#1}{#2}}
\newcommand{\f}[1]{\texttt{\detokenize{#1}}}
\newcommand{\pg}[1]{\clearpage\section{#1}}
\setlength{\emergencystretch}{8em}
\begin{document}
\section{人选中一件东西，AI 就有了上下文}
你在场景里点选龙，拖动粗糙度滑块。引擎知道你说的是哪条龙、这条龙用了哪个模型、当前编辑的是第几个材质槽。现在把手从鼠标移开，说一句「表面再粗糙一点，颜色保留」。对于一款面向人的 AI 原生编辑器，这句话应当接得上刚才的操作。改完你可以接着拖同一个滑块，可以撤销，也可以关掉聊天后保存项目。这个例子决定了架构中最重要的边界：谁有资格解释当前场景，谁有资格修改它。

Redflag 的属性面板已经掌握这份上下文。\f{SceneHierachyPanel::DrawComponents} 拿到选中的实体。在模型材质区域，面板能显示子网格原材质，启用实例覆盖时会克隆材质，然后用 \f{ImGui::SliderFloat} 直接写入粗糙度。这个写法对一个面板很自然。现在有了第二个调用者，就值得把「修改该实体的覆盖材质」提成一个普通编辑器函数。面板继续决定滑块怎么画，聊天助手继续理解「再粗糙一点」的意思，两个入口都把确定的目标和值交给这个函数。

{\small\color{muted}现有入口：\f{Redflag/src/Redflag/Scene/Panel/SceneHierachyPanel.cpp}，\f{DrawComponents}，模型覆盖材质滑块见第 1667 行。下文新增的接口及界面均属拟议设计。}

\begin{center}
\begin{tikzpicture}[n/.style={draw=teal!65,fill=teal!5,rounded corners=2pt,align=center,minimum height=13mm,font=\small,text width=31mm},a/.style={-{Stealth[length=2mm]},thick,teal}]
\node[n](human) at (0,0){属性面板\\接收鼠标操作};
\node[n](agent) at (0,-1.9){聊天助手\\读取选择与工具说明};
\node[n](edit) at (5,-.95){编辑器主线程\\执行编辑，登记撤销};
\node[n](scene) at (10,-.95){编辑场景\\拥有实体和组件};
\draw[a](human.east)--(edit.west);\draw[a](agent.east)--(edit.west);\draw[a](edit.east)--(scene.west);
\end{tikzpicture}

{\small 图 1\quad 两个入口提交同一项编辑。箭头表示谁把操作交给谁。}
\end{center}

已有引擎仍按 Application、Layer、Scene、Entity 组织。AI 工具放在制作侧；发布出去的游戏读取普通场景、资产和脚本。人的层级面板、场景视图、属性面板和项目资源仍是四个主要入口。任务结果可以直接定位到这些位置，例如点击「龙的材质已修改」就选中龙和对应的槽位。一句话生成的大项目也应当能这样逐项接手。

\textbf{设计出处。} Amershi 等在 CHI 2019 的人机交互准则中讨论了让用户方便地调用、纠正和取消 AI 帮助。Redflag 可以把它落实成三个可测操作：从当前选择发起任务、手工修正任一生成对象、用同一份编辑历史撤销。 \src{https://www.microsoft.com/en-us/research/publication/guidelines-for-human-ai-interaction/}{Guidelines for Human-AI Interaction}

\pg{一次材质修改怎样安全落到场景里}
助手先查「当前选中了什么」。编辑器返回当前文档、编辑或试玩模式、实体 UUID、材质槽和实际值。人读到「龙」，工具使用明确的标识。已有参考场景中龙的实体号为 1，模型号却是 6122188709592601978；两者描述的是放在场景里的这一条龙和磁盘上可复用的模型。向 JSON 客户端传输时把二者编码为不同字段的十进制字符串，以保留 64 位精度。查询结果只附本次可操作的字段，分量范围与单位紧跟字段。需要完整子网格列表时再查询。

假设槽 0 已有实例覆盖，读出的粗糙度约为 0.63。AI 可以选 0.80 作为「更粗糙」的试调值。主线程收到请求时重新按 UUID 取实体，读取槽位和当前值。若你在它思考时把同一数值改成了 0.74，编辑器返回当前值和冲突字段，助手重新判断。一盏灯在同一时间被改名，无须阻止改龙。对于刚被重新导入的模型，还要核对槽位对应的子网格，避免把皮肤的材质写到衣服上。

\begin{center}\small
\begin{tabularx}{\linewidth}{@{}p{28mm}XX@{}}\toprule
拟议调用 & 由谁提供 & 调用者下一步能做什么\\\midrule
检查选择 & 编辑器返回实体、所在文档、槽位、当前字段、取值范围、是否已有覆盖 & AI 解释目标，面板能定位同一对象。\\
修改材质 & 两个入口提交目标、字段、新值和本次编辑所依据的旧状态 & 主线程返回实际变更或冲突；渲染读取变更后的材质。\\
拍摄对照图 & 编辑器按给定机位、尺寸和曝光绘制，返回图像和对应场景版本 & 用户可并排比较；AI 可检查是否动错对象。\\
撤销操作 & 编辑历史保存该操作的前后值与覆盖状态 & 原来继承模型材质时，撤销后应恢复继承。\\
\bottomrule\end{tabularx}\end{center}

手拖滑块时，按下到松开产生一次撤销记录，中途仍能看见连续预览。AI 的一次操作可以包含一组明确的改动，先检查所涉及的目标，再统一应用。用户若随后又手调，应按编辑历史顺序撤销；单独撤回旧任务则要检查它与后续操作是否修改了同一字段。价格较高的生成任务可以取消，已拿到的候选素材仍可查阅。

Unity 的 SerializedObject 文档里，同一对象的两个编辑数据流需要重新读取变化并显式提交。这里建议按本次目标字段核对旧状态，避免聊天中稍早的预览覆盖刚刚完成的人工调整。 \src{https://docs.unity3d.com/6000.0/Documentation/ScriptReference/SerializedObject.html}{Unity 6，SerializedObject}

\f{ForgeLayer::OnScenePlay} 用 \f{Scene::Copy} 建立试玩场景，\f{SaveScene} 保存编辑场景。两者已区分，但 \f{Scene::Copy} 按组件复制，其中的材质引用可以指向同一个对象。因此试调外观时，新建覆盖材质后再写值，或者用独立进程打开项目副本。验收要同时检查原场景、其他实体和试玩结束后的状态。 \src{https://docs.godotengine.org/en/stable/classes/class_editorundoredomanager.html}{Godot 编辑器操作历史，可对照研究}

\pg{聊天助手怎样调用引擎}
以「给这个房间加一个出口，保留我摆好的主相机」为例。聊天适配器先把当前选区、场景、可用资产和游戏脚本 API 的相关部分交给模型。模型可以询问哪一面墙空着，读取包围盒，提出一个出口位置。引擎提供有明确输入输出的小操作，由它负责解析 ID、验证位置并写入对象。模型每次看到的是这一步真实返回的数据。房间布局、资源身份和编辑历史由项目保存，换一个聊天仍可重新读出。

个人项目只需一个本地助手和一个制作工具入口。脚本或 CLI 适合编译、下载、导入和打包。正在运行的编辑器要读未保存的场景，适合经本机连接把请求排到主线程，在下一次安全的更新点执行；慢的网络生成和编译在独立进程执行，只回传结果。MCP 可以描述工具的名称、参数和返回内容。把一个 MCP 操作映射到已有编辑器调用即可，同一操作也可以从编辑器菜单运行。游戏每帧的输入、动画和渲染走原来的运行时路径。

工具说明本身就是接口的一部分。以读取场景为例，可先返回实体名、UUID、组件名和包围盒，查询一个实体时再附材质或动画数据。失败要给出下一步能使用的信息，例如「当前是试玩副本，以下是可编辑的场景对象」或「模型有两个材质槽，此请求引用槽 2」。长日志、帧捕获和图片留在本地，每次先返回摘要和产物位置。有需求再加新操作，维持个人能读懂的工具数量。 \src{https://arxiv.org/html/2405.15793v3}{SWE-agent，§3 与 §5.1}\quad \src{https://www.anthropic.com/engineering/writing-tools-for-agents}{Writing effective tools for AI agents}

\textbf{MCP 和 skill 分别放什么。} Maya 桥负责执行具体的骨架查询与烘焙，返回关节、帧和导出的文件。项目的「角色导入」skill 写清该先查哪些关节、如何对齐单位、用哪些动作验收。角色最终仍作为可查看的模型、骨架和动作保存在项目里。已有 RenderDoc MCP 同理：工具能打开捕获，项目 skill 说明在 Redflag 里该找哪个 Pass、保存哪些附件、如何用同输入复测。工具可替换，已导入的角色和场景不随助手变动。

\textbf{生成过程可以中断。} 一个本地任务目录记录用户需求、输入文件版本、外部任务号、已产生文件和通过的检查。程序恢复时先查已有任务号和文件哈希，避免重复付费生成。把候选资源写入新目录，校验通过后登记到项目。遇到风格选择，可以按你的默认偏好直接采用并保留候选；有意挑选时，编辑器展示来源、许可和实拍。 \src{https://www.anthropic.com/research/effective-harnesses-for-long-running-agents}{长任务状态与交接}

\pg{一句话怎样变成一个可编辑的游戏}
假设你说「做个小仓库，两个人抢能量块，先得三分获胜」。助手先把它解释成玩家能体验到的几件事：两人从不同位置出现、靠近能量块可以拾取、双方看到相同的分数、三分弹出结算、可以重新开始。默认地图尺寸、操作键和美术风格取自你的个人项目设置。这份简短说明放在项目旁边，供你日后修改。明确常用缺省值后，完整生成可以从一句话开始。

第一批产物是普通场景对象。主相机、出生点、能量块都在层级中有名字和位置。GPT Image 的图提供构图与颜色，资源查询提供实际尺寸；布局工具依据接地、碰撞、间距和通路计算变换。你可以圈选一处说「只把这里改成侧门」。已有实体按 ID 定位，生成差异限定到这处修改，重开场景还能看见你亲手放的灯。Holodeck 让模型提供空间关系，再用约束布局，这一分工可借到 Redflag。 \src{https://openaccess.thecvf.com/content/CVPR2024/html/Yang_Holodeck_Language_Guided_Generation_of_3D_Embodied_AI_Environments_CVPR_2024_paper.html}{Holodeck，CVPR 2024}

玩法规则落在游戏自己的 C\# 源文件里。AI 查询当前版本可用的组件、方法和示例，写脚本后调用普通构建流程，读编译器返回的文件和行号。像「胜利分数」这样你常调的量应该是可编辑字段。你可以在属性面板把三分改成五分；稍后要求 AI 换皮肤时，它应只改资源及引用。游戏 UI 的比分绑定到游戏状态。更改文字位置用布局属性，更改得分条件打开规则脚本。动画状态机也一样，应能在属性中查看当前状态、触发条件及所选片段。

当前还有一个实在的构建缺口。\f{SandboxProject/premake5.lua} 将游戏编成固定的 \f{Sandbox.dll}，\f{ScriptEngine::LoadAssemblies} 按固定位置寻找它。建议在可编辑的项目设置中保存启动场景、资源根目录和游戏程序集。编辑器点击 Play、AI 试玩、独立播放器启动时读取同一份设置，才能用第二个生成项目验证它没有误加载第一个游戏的脚本。

从生成目标到文件落地，中间存在两种工作尺度。新建游戏时可以在新项目目录一次生成场景、脚本、UI 和资产清单，构建成功后打开项目。继续修改时，先读磁盘版本与当前打开的未保存场景，按对象和字段生成差异。跨场景与脚本的修改在独立项目副本中编译、试玩；采用时重新检查基准文件，冲突直接定位到实体或源码。小型滑块编辑使用进程内撤销，大范围生成修改用项目版本或文件差异还原。这样一次编译失败也不妨碍你继续使用原项目。

\textbf{边界由实际需求决定。} 若编译器报告当前脚本 API 没有拾取事件，助手先检查能否用已有的触发器或距离查询表达。已有能力足够时把逻辑留在游戏。确实需要新的底层操作时，注明缺少的引擎接口、调用者和验收条件，作为单独的引擎变更。终版可以自动完成已支持范围的游戏，新的玩法也保留普通代码扩展入口。

\pg{从属性面板进入源码，再回到同一张场景}
你正在看能量块，想把它改成「连续拾取两次才加分」。在拟议的编辑器中，选中它就能看见挂载的脚本类型、源文件和本次运行加载的程序集。点「打开源码」可以亲手修改，也可以选中这个对象对 AI 说「改拾取逻辑」。两条路径都编辑项目里的普通 C\# 文件。场景只保存脚本类型、实体引用和序列化参数。以后你在 IDE 里加的提前返回、注释和新分支，都成为 AI 下次读取的真实输入。聊天中最初生成它的那份文字说明只用来理解当时意图。

Redflag 已有从面板打开脚本的 \f{OpenScriptSourceFile} 和 \f{RecompileScriptAssemblies}，都在 \f{SceneHierachyPanel.cpp}。现有重编入口检查是否正在 Play，停止运行后卸载程序集，调用 MSBuild，再调用 \f{ScriptEngine::LoadAssemblies}。这个入口目前固定查找 SandboxProject。拟议做法是让它从当前项目设置取得工程与程序集路径，保留你熟悉的手动编译按钮，并给 AI 返回同一份编译诊断、源码位置和加载结果。启动回放前记录程序集的内容哈希，避免把旧 DLL 的成功当作新代码已经生效。

\textbf{AI 怎样与你同时改文件。} 开始任务时，先读工作目录的实际文件，包括尚未提交的修改与任务依赖的未跟踪文件。可用 Git diff 辅助找范围，单从 HEAD 新建的工作树并不包含这些内容。小改动按读取时的文本生成局部补丁；跨文件或引擎修改在独立工作目录构建，用这批文件的实际基准初始化。应用前比较基准、候选和当前内容，能合并的不相关改动直接合并，同一处冲突保留你的版本和候选差异。未落盘的 IDE 缓冲区只有接入 IDE 后才能读取；普通本地工具遇到这类文件时保留候选补丁，等保存后再合并。场景里的未保存编辑则向正在运行的编辑器查询。两种未保存状态由各自的拥有者解释。

\textbf{随时编辑，按对象选择生效时机。} 游戏参数在编辑场景中立即预览。C\# 更改按现有约定停 Play、重编、重载后重新运行；运行中若希望保留关卡进度，需为该游戏明确实现状态迁移。Shader 源文件可随时修改，拟议的热重载在新程序编译、链接且资源绑定兼容后，于帧边界切换；绑定布局变化时重建依赖。修改引擎 C++ 类型、内存布局或资源所有权后，采用独立输出目录构建新编辑器，保存当前场景或恢复副本，再用新进程打开它。构建失败保留诊断与源码，当前可用二进制继续留存。

\textbf{一个可验证的接力。} 让 AI 生成拾取脚本，你增加一条「结算后禁止得分」，再让 AI 增加双倍积分区域。构建后回放三次拾取：普通区域得一分，双倍区域得两分，结算后得零分。最后从属性面板打开文件，核对你写的规则和改动差异。这个用例检验的是人和 AI 能接续开发同一份源码。

\pg{当一句需求需要改引擎源码}
现在说「让龙的拉丝材质有一个可以旋转的高光方向」。有了粗糙度并不等于有各向异性的方向参数。助手先查当前材质字段、场景读写、实际绘制与着色器，判断需求能否由已有能力表达。若本项目确实要支持这一表面模型，AI 可以像一名引擎开发者那样修改 C++ 和 Shader，连同界面与文件兼容性一起交付。你也可以亲自写着色计算，让 AI 只负责接编辑器和验证。

从人的操作倒推，属性面板里需要选中龙、打开实例材质、转动方向，在当前灯光下看高光，再保存重开。工程上先约定方向相对哪条切线、角度单位及哪些网格具备有效切线，再定义各向异性强度。方向旋转只在强度非零时有可观察的作用；强度取零应能复现旧效果。原文件没有这些字段时采用该默认值。若实际 Shader 缺必要切线信息，工具要报告具体输入限制。这里描述的是拟议特性，没有为它给出未经验证的渲染公式或性能结论。

\begin{center}\small
\begin{tabularx}{\linewidth}{@{}p{32mm}XX@{}}\toprule
源码改动落点 & 调用关系 & 必须核对的结果\\\midrule
材质与持久化 & \f{Material.h} 定义字段，属性面板编辑实例，\f{WriteMaterial} / \f{ReadMaterial} 保存恢复 & 旧场景仍能加载；新值重开一致；同模型另一实体保留原状。\\
渲染与效果 & 沿活跃后端追踪场景选材质、参数上传、片元阶段读取；OpenGL 网格着色源码入口为 \f{Renderer3D_Mesh.glsl} & 固定机位、几何、曝光和灯，比较零强度与两个方向，标明已验证的后端。\\
工具与手册 & 从同一编辑操作提供字段含义、范围、单位和可用条件，更新脚本绑定或工具描述中实际向外暴露的部分 & 人拖滑块和 AI 写同一实参得到一致状态；旧工具若不认识该字段可查询新能力。\\
\bottomrule\end{tabularx}\end{center}

\textbf{从引擎变更回到制作任务。} 编译通过后先用新构建打开测试场景，运行旧材质兼容用例与新增对照，保留图像、程序版本和失败项。通过后，工具重新查询新编辑器支持的字段，再对原任务中指定的龙应用材质修改。交付时你能分别看到哪几行扩大了引擎能力、哪一项更改只属于龙，以及上次能运行的程序在哪里。以后在 IDE 里手改 C++，同样从构建入口刷新能力和验证结果。

不需要给每个 C++ 函数制作一个 MCP 工具。对稳定制作动作，工具调用已经编译好的编辑器操作；需求超出这些动作时，使用普通代码编辑、编译器和测试来扩展项目。两条路径汇合在同一套源码、同一个编辑器和同一张可编辑场景上。

\pg{AI 怎样知道交付的游戏能玩}
「能量块被拾取」不是一个屏幕像素。它涉及输入、角色位置、玩法状态和 UI。让 AI 自己试玩时，测试进程应走正常玩家入口，输入一段移动指令，等到游戏报告拾取事件，再读取场景状态、双方比分与结算画面。读状态是为了确定规则对不对，看画面是为了确定玩家能不能看懂。断言不能直接设置比分，否则会跳过真正需要验证的交互。

以最后一分为例，记录服务器认定的得分事件、两端显示的比分、结算控件和重开后的初始值。再用另一条输入回放让玩家经过目标旁边却没有碰到它，预期比分维持原值。保留一条并未交给生成助手的验收用例，避免它只为成功演示写一条固定路径。最终从干净的发布目录再启动一遍，验证入口、资源引用和程序集。人直接拿到可游玩的构建，同时有可改的源项目、操作说明和已知限制。 \src{https://arxiv.org/abs/2602.11103}{GameDevBench}\quad \src{https://www.anthropic.com/research/demystifying-evals-for-ai-agents}{代理结果评测}

还有一个今天就能从源码核对的反例。属性面板能替换覆盖材质的贴图，但当前 \f{SceneSerializer.cpp} 中的 \f{WriteMaterial} 和 \f{ReadMaterial} 只处理底色、粗糙度、透明模式等参数，没有读写覆盖贴图引用。即使预览里换对了图，也要保存后用新场景实例重新打开，检查同一槽的资源引用与实拍。工具应分别报告「预览通过」和「重新打开通过」；缺持久化能力时指出这个真实限制。

当你说「做出的房间一转头就糊」，助手应先固定镜头、分辨率、灯光和输入，并记录连续帧。它可以读取引擎已有的时间线和图像误差，比较开关效果后的差异，必要时用 RenderDoc 找到第一个输出异常的事件。修改后用相同回放复测。如果需求是提高帧率，则另外跑相同构建配置与预热条件的计时，分别报告 CPU、GPU 及画质。技术依据、假设与仍待人判断的画面风格，可以在同一次交付中各自说清楚。

\textbf{衡量这套架构有没有帮到你。} 留下两种结构不同的小游戏需求，重复测试四件事：一句话能否生成并启动；你能否在编辑器里独立找到并修改胜利条件；你改完后，AI 能否继续加内容并保留这项手工修改；任务失败后能否恢复到最后一个可玩的版本。记录成功次数、你实际插手的步骤、时间和生成费用。对个人引擎而言，能自己接着改，与生成成功同样是产品验收条件。

\textbf{阅读顺序。} 做编辑器时先读 Amershi 等关于用户纠正与控制的建议，再看 Godot 如何给一次编辑配上执行和撤销动作。做 AI 工具时读 SWE-agent 的接口实验和 Anthropic 的工具设计文章。设计场景生成读 Holodeck；设计自动试玩读 GameDevBench。这些工作分别提供交互、接口、空间布局和评测的依据。本文的 Redflag 编辑调用、任务目录和验收用例均为针对本项目的设计。
\end{document}

